\section{Fragestellung: dasselbe Objekt --- was HLV damit macht, bleibt HLVs Sache}
Eine Aussage steht am Anfang und soll unmissverständlich sein: FFGFT und HLV arbeiten mit
\emph{demselben} geometrischen Objekt. Das ist keine Deutung und keine These, sondern eine
nachrechenbare Tatsache, die unten belegt wird -- das $\varphi$-ikosaedrische Achsensystem, auf
dem $\theta=2/9$ als Invariante sitzt. Alles Weitere in diesem Dokument ist Folgerung aus dieser
einen Feststellung.

Davon strikt zu trennen ist eine zweite Frage, die hier \emph{nicht} gestellt und \emph{nicht}
beantwortet wird: was HLV mit diesem Objekt macht -- welche Vortex- oder Randmoden-Deutungen,
welche Ableitungen und welche Formulierungen HLV darauf aufbaut. Das ist HLVs Sache, nicht die
der FFGFT. Ob diese Formulierungen korrekt sind, wird hier weder behauptet noch geprüft noch
adressiert. Die FFGFT stellt das geteilte Objekt fest und -- über $E_0$ -- eine Prüfrichtung
bereit; für den Gebrauch, den HLV vom Objekt macht, ist HLV allein verantwortlich.

Der Zweck des Dokuments ist damit bewusst zweigeteilt und bescheiden. Es hält fest, dass beide
über demselben Objekt rechnen, und dass FFGFT eine Möglichkeit bietet, HLV an der Empirie zu
messen; und es zieht zugleich die Grenze, jenseits derer über HLV nichts ausgesagt wird.

\section{Das geteilte $\varphi$-ikosaedrische Objekt}
In der verfügbaren rechnerischen Darstellung von HLV -- dem Cut-and-Project des
Homologie-Prüfskripts \texttt{ffgft\_hlv\_homology\_check\_v3.py} -- wird der physikalische
Raum durch den ikosaedrischen Projektor mit dem Parameter $\tau=\varphi$ erzeugt,
$\varphi=\tfrac{1+\sqrt5}{2}$. Dessen sechs projizierte Achsen sind die sechs fünfzähligen
Achsen des Ikosaeders: alle fünfzehn paarweisen Winkel sind gleich $\arccos(1/\sqrt5)=
63{,}435^\circ$. Genau auf diesen Achsen lebt die $C_3$-in-$A_5$-Verdopplung der FFGFT
(Dok.~285, Dok.~293).

Es handelt sich also nicht um zwei ähnliche Strukturen, sondern um dasselbe Objekt. Wendet man
auf diese Achsen die fünfzählige Umverteilung der $C_3$-Moden an (Dok.~293), so ist der Anteil
in die triviale Mode exakt $p_0=2/9$. Der Koide-Winkel $\theta=2/9$ ist damit eine
\emph{Invariante auf dem geteilten Objekt}, nicht eine Übersetzung des einen Rahmens in den
anderen. Der Wert ist überdies spezifisch für den Goldenen Schnitt: nur $\tau=\varphi$ macht
das Winkelspektrum exakt uniform und den trivialen Anteil exakt $2/9$; der silberne Wert
$\tau=1+\sqrt2$ ergibt $0{,}2534$ und ein bereits messbar verschobenes Spektrum (Winkel
$45^\circ$ und $69{,}3^\circ$). Das geteilte
Objekt ist also nicht irgendein Cut-and-Project, sondern das $\varphi$-ikosaedrische.

Diese Aussage ist rechnerisch und ohne Zirkel: sie folgt aus dem Vergleich der Projektor-Achsen
und der Modenumverteilung, nachprüfbar in \texttt{ikosaeder\_theta\_delta.py} (Seed
$20780458$) und in der reproduzierbaren Winkelspektrum-Prüfung \texttt{a5\_bridge\_discrimination.py}.

\section{Abgrenzung: welches HLV-Objekt geteilt ist}
Das hier festgestellte geteilte Objekt ist die \emph{geometrische} Schicht von HLV -- das
$\varphi$-ikosaedrische Achsensystem, auf dem $\theta=2/9$ als Invariante sitzt. Davon strikt zu
trennen ist ein zweites HLV-Objekt gleichen Namens: die \emph{temporal-dynamische} Spiral-Zeit, ein
Nicht-Markov-Operator mit triadischer Zeitkoordinate $\Psi(t)=(t,\phi(t),\chi(t))$ und
Gedächtniskanal (Zenodo 20643462). Diese temporal-dynamische Spiral-Zeit trägt weder $A_5$ noch den
ikosaedrischen Projektor; sie ist \emph{nicht} das hier geteilte Objekt und fällt nicht unter die
vorliegende Feststellung.

Die Prüfrichtung dieses Dokuments betrifft allein die geometrische Schicht. Der Abgleich der
temporal-dynamischen Spiral-Zeit mit FFGFTs Zeit-Windung wird gesondert geführt (Dok.~295); dort
zeigt eine rechnerische Prüfung, dass jenes Objekt das Selbstähnlichkeits-Verhältnis $e$ der
FFGFT-Log-Spirale nicht trägt -- auf seinen eigenen Definitionen ist es eine Windung konstanter
Ganghöhe, nicht äquiangular. Beide Abgrenzungen zusammen halten die Rede von \enquote{geteilt}
sauber: geteilt ist das $\varphi$-ikosaedrische geometrische Objekt, nicht HLVs Spiral-Zeit als
Zeitoperator. Die geteilte Invariante zwischen den beiden Schichten ist allein die $Z_3$
($C_3<A_5$); die Dimensionszahlen der beiden Rahmen ($3$, $6$) decken sich, ihre Bedeutung nicht
(Dok.~295).

\section{Die $E_0$-Asymmetrie}
Das geteilte Objekt ist skalenlos. Die Geometrie kennt weder Masse noch Energie; sie liefert
reine Zahlen -- den Winkel und die Nenner-$9$-Gewichte (Dok.~293). Hier setzt die eigentliche
Asymmetrie zwischen den beiden Rahmen ein.

FFGFT verfügt über einen Referenzpunkt $E_0$ (Dok.~292), der das skalenlose Muster an die
gemessenen Leptonen $e,\mu,\tau$ bindet und die absolute Lage fixiert. Damit kann FFGFT vom
Muster \emph{empirisch zurückrechnen}: aus dem geteilten Objekt, den Nenner-$9$-Gewichten und
$E_0$ folgen absolute Massen. HLV besitzt keine solche Skala. Es trägt das Muster -- das
geteilte geometrische Objekt --, aber keinen Anker in Energie oder Masse. HLV kann daher aus
dem Objekt allein keine Massen gewinnen. Es kann es nur, wenn es $E_0$ von FFGFT übernimmt; in
diesem Fall aber rechnet die Empirie der FFGFT, lediglich in den Koordinaten von HLV
ausgedrückt.

\section{Die Prüfrichtung: von FFGFT nach HLV}
Aus der geteilten Geometrie und der einseitigen Skala folgt eine gerichtete Beziehung. Weil
FFGFT die fehlende Skala beisteuert, bietet FFGFT eine Möglichkeit, den geometrischen Gehalt
von HLV an den Daten zu prüfen: man versieht das geteilte Objekt mit $E_0$ und testet, ob es
die Leptonmassen reproduziert. Das ist eine Verifikationsrichtung von FFGFT nach HLV.

Es ist ausdrücklich \emph{keine} Bestätigung, dass HLVs eigene Formulierungen korrekt sind. Das
geteilte Objekt besagt allein, dass beide dieselbe Geometrie benutzen -- nicht, dass die
Interpretationen und Ableitungen, die HLV auf dieser Geometrie errichtet, gültig sind. Die
Prüfrichtung testet die Entsprechung von Geometrie und Massen mit den Mitteln der FFGFT; sie
berührt HLVs eigene Konstruktion nicht und adelt sie nicht.

Umgekehrt gilt: HLV liefert der FFGFT empirisch nichts, was diese nicht schon hätte. Das Muster
ist geteilt, die Skala hat FFGFT selbst. Es gibt keinen Empirie-Transfer von HLV nach FFGFT,
und FFGFT braucht HLV für die Massen nicht. Der empirische Pfeil ist einseitig.

\section{Was der Befund besagt und was nicht}
Belegt und rechnerisch gesichert ist: das geteilte $\varphi$-ikosaedrische Objekt; $\theta=2/9$
als dessen Invariante; die Spezifität für $\tau=\varphi$; und die $E_0$-Asymmetrie, aus der die
Prüfrichtung folgt.

Nicht belegt und nicht behauptet ist: die Korrektheit der Formulierungen, die HLV enthält; eine
unabhängige empirische Bestätigung der FFGFT durch HLV; und die Frage, ob der ikosaedrische
Projektor der fundamentalen Behauptung von HLV \emph{intrinsisch} ist oder in den
Prüfskripten eine Modellierungswahl darstellt. Letzteres ist von HLV selbst zu klären und wird
hier offengelassen.

\section{Methodische Einordnung}
In den drei methodischen Schichten der FFGFT ordnet sich der Befund klar ein. Das geteilte
Objekt ist ein Brücken-Befund: eine geometrische, algebraisch belegbare Entsprechung. Die
Zurückrechnung auf Massen ist Kern-Ableitung, sie folgt aus $\xi$ und $E_0$ und ist
FFGFT-intern. Der Beitrag von HLV zur Empirie ist null. Im Sinne von P35 bleibt die Aussage,
HLV sei korrekt, ein Kandidat und keine Ableitung; abgeleitet ist allein das geteilte Objekt.

\section{Der erste Prüfversuch ist überholt}
Der erste Prüfversuch -- der Homologie-Check \texttt{ffgft\_hlv\_homology\_check\_v3.py}, ein
grober topologischer Punktwolken-Diagnostiker -- ist durch die exakte Hilbertraum-Variante
(Dok.~293) überholt und war für die tragende Frage ohnehin strukturell blind: die persistente
$H_1$ einer Punktwolke sieht die Achsen-Winkel-Struktur nicht, an der $2/9$ hängt, sodass
$\tau=\varphi$ (exakt $2/9$) und der Kontrollwert $\tau=\sqrt2$ ($0{,}2139$) topologisch nahezu
ununterscheidbare Wolken erzeugen. Sein früherer Befund einer \enquote{nur schwachen Trennung
von generisch} ist deshalb keine Aussage über $\varphi$ oder $A_5$, sondern die erwartbare
Blindheit des groben Statistikums. In der spektralen Variante erscheint $\theta=2/9$ dagegen
unmittelbar als exakte Invariante $p_0=\lvert\langle v_0|R_5 v_e\rangle\rvert^2=2/9$; der Befund
vom geteilten Objekt ruht auf ihr, nicht auf dem Homologie-Diagnostiker.

\section{Der zweite Test: was er auflöst und wo er zirkulär wird}
Der zweite Prüfversuch -- das Instrument \texttt{a5\_bridge\_discrimination.py}
(Seed $20780458$) -- verwendet genau die Statistik, die dem ersten fehlte: das Winkelspektrum
der sechs projizierten Achsen, gemessen als Abstand zum ikosaedrischen Referenzspektrum, in dem
alle fünfzehn paarweisen Winkel gleich $63{,}435^\circ$ sind. Damit löst er die Achsen-Winkel-
Struktur auf, an der $2/9$ hängt. Er ist so kalibriert, dass ein exakt ikosaedrisches Objekt den
Abstand null ergibt, während generische zufällige Tight-Frames und eine echt generische
$\sqrt2$-Konstruktion in einem breiten generischen Bereich liegen. Er ist damit das Instrument,
das $\varphi$ und $A_5$ sehen kann, wo der erste blind war.

Was der zweite Test auf der verfügbaren Darstellung erbringt, ist zweierlei. Erstens liest der
HLV-Projektor mit $\tau=\varphi$ exakt ikosaedrisch: alle fünfzehn Winkel $63{,}435^\circ$,
Abstand null, und die fünfzählige Umverteilung auf seinen Achsen ergibt exakt $2/9$. Das ist
jedoch eine Folge der Konstruktion -- das Objekt \emph{ist} der ikosaedrische Projektor --, die
Lesung also zirkulär und keine unabhängige Bestätigung von $A_5$. Das Instrument behandelt dies
korrekt nicht als Urteil, sondern macht die Zirkularität ausdrücklich. Zweitens weist es
quantitativ nach, dass die frühere \enquote{generische} Kontrolle mit $\tau=\sqrt2$ nicht
generisch ist: RMS-Abstand $3{,}47$ zum ikosaedrischen Referenzspektrum (das exakt
ikosaedrische Objekt hat Abstand $0$), ein zweiwertiges Spektrum $61{,}87^\circ/70{,}53^\circ$, und
ein Umverteilungsgewicht $0{,}2139$, das nicht $2/9$ ist. Dieser Abstand liegt weit unterhalb des
generischen Bandes zufälliger Frames (Median-RMS etwa $22$, $5$--$95\,\%$ rund $17$--$28$): die
$\sqrt2$-Kontrolle ist also strukturiert und nicht generisch. Damit ist die $\varphi$-Spezifität
festgenagelt: allein $\tau=\varphi$ macht das Spektrum exakt uniform und das Gewicht exakt
$2/9$.

Der Ertrag des zweiten Tests ist somit diagnostisch und instrumentell, nicht bestätigend. Er
stellt die richtige, reproduzierbare Statistik bereit, die auflöst, was der erste nicht auflösen
konnte; er quantifiziert die Zirkularität einer Prüfung des verfügbaren Objekts und die
Kontamination der früheren Kontrolle; und er isoliert $\varphi$ als den ausgezeichneten Wert.
Was er noch nicht liefert, ist eine unabhängige Bestätigung: die verlangte ein HLV-Objekt, das
\emph{ohne} eingebackenen ikosaedrischen Projektor dargestellt ist. Bis ein solches Objekt
vorliegt, steht das Instrument bereit, und der Befund vom geteilten Objekt ruht
auf der exakten spektralen Invariante, nicht auf einer nicht-zirkulären Diskrimination.

\section{Status}
Gesichert ist das geteilte $\varphi$-ikosaedrische Objekt und die $E_0$-Asymmetrie, und daraus
eine Prüfrichtung von FFGFT nach HLV: FFGFT bietet die Möglichkeit, den geometrischen Gehalt von
HLV an den gemessenen Massen zu prüfen. Nicht behauptet wird irgendeine Korrektheit der
Formulierungen, die HLV enthält. Offen bleibt ein Punkt: ob der ikosaedrische Projektor HLV
intrinsisch ist. An ihn gebunden ist die einzige noch mögliche unabhängige Prüfung -- eine
nicht-zirkuläre Diskrimination, die ein HLV-Objekt ohne eingebackenen ikosaedrischen Projektor
verlangt; das Instrument \texttt{a5\_bridge\_discrimination.py} steht dafür bereit.
Der frühere Homologie-Kontrast ist damit erledigt, nicht offen.

$\theta=2/9$ selbst steht unabhängig von all dem. Es folgt aus der $Z_3$-Zirkulant-Diagonali\-
sierung und der ikosaedrischen Umverteilung (Dok.~293) und braucht HLV in keiner Weise. Die hier
festgehaltene Beziehung ist eine geteilte Geometrie mit einseitiger Prüfrichtung, nicht eine
gegenseitige Bestätigung.
